\section{Two exact matrix backends}
\label{sec:backend}

We give two versions of the result. The first uses the classical field
bounds already available to the incoming rectangular packet. The second
also imports the stronger published 2026 matrix-multiplication theorems.
Keeping both versions makes the additional dependency explicit.

\subsection{A common rectangular formula}

Suppose fixed bilinear matrix algorithms are available with dual aspect
ratio $a\in(0,1)$ and square exponent $w>2$, allowing arbitrarily small
positive exponent slack. For sorted base-two logarithmic dimensions
$x\ge y\ge z\ge0$, define
\begin{equation}\label{eq:rectcost}
 B_{a,w}(x,y,z)=\max\left\{x+y,\quad
      x+\frac{1-a(w-1)}{1-a}y+\frac{w-2}{1-a}z\right\}.
\end{equation}
For every fixed $\delta>0$, there is a fixed exact algorithm with time a
polynomial factor times $2^{(1+\delta)B_{a,w}}$. Its space is a polynomial
factor times the largest input or output array.

Indeed, block the largest dimension at scale $2^y$. For $z\le ay$ use
the dual-exponent algorithm. Otherwise tensor the dual and square schemes
at logarithmic scales
\[
 t_0=(y-z)/(1-a),\qquad t_1=(z-ay)/(1-a).
\]
They satisfy $t_0+t_1=y$ and $at_0+t_1=z$. Their rank exponent, plus
the blocking cost, is exactly the second branch of \eqref{eq:rectcost}.
Fixed-power padding changes sizes by constant factors. Choosing each
component's exponent slack sufficiently smaller than $2\delta$ gives the
stated multiplicative allowance, since $B_{a,w}\ge x+y\ge2y$.
The first branch is the essential input/output cost.

For the classical version use $a=3/10$ and $w=19/8$, justified by the
strictly stronger results of Vassilevska Williams, Xu, Xu and Zhou
\cite{mmclassical}. Then
\begin{equation}\label{eq:classicalbackend}
 B_{\rm cl}=\max\{x+y,\ x+(47/56)y+(30/56)z\}.
\end{equation}
For the additional 2026 version use $a=9/20$ and $w=9/4$.
The published results give $\omega(\mathbb C)\le9/4$ and a
characteristic-zero dual exponent greater than $0.465$, hence greater than
$9/20$ \cite{mmnew,mmrect}. Therefore
\begin{equation}\label{eq:newbackend}
 B_{\rm new}=\max\{x+y,\ x+(35/44)y+(20/44)z\}.
\end{equation}
We fix $\delta=10^{-5}$ in the final schedules. This charges the
positive slack of the algorithms explicitly. The endpoint statement
$\omega\le9/4$ by itself is not being interpreted as one fixed algorithm
of exact cost $N^{9/4+o(1)}$.

\subsection{Transfer from complex arithmetic to exact bit arithmetic}

The stronger result is stated over $\mathbb C$, whereas the SAT contractor
uses rings with zero divisors. The transfer requires an argument.

Fix the positive tolerance once. The matrix theorem supplies a finite
exact rank scheme. Its polynomial equations have rational coefficients;
if they have a complex solution, they have a solution over
$\overline{\Q}$. All coefficients of this fixed scheme lie in a single
number field $K/\Q$ of fixed degree. The square paper explicitly records
this observation, and the rectangular paper states its
characteristic-zero transfer \cite{mmnew,mmrect}. Choose the finitely
many schemes used by the two tensor factors in a common number field.

Represent elements of $K$ in a fixed rational basis. A field operation
becomes a fixed number of rational-coordinate operations. Recursive
bilinear schemes form only fixed linear combinations of input blocks
until their leaf products, then fixed linear combinations of outputs.
Consequently their intermediate coordinate bit lengths are polynomial in
the input bit length and the logarithm of the matrix dimensions. Fixed
denominators and coefficient heights accumulate over only logarithmically
many recursive levels. This gives a uniform exact rational implementation
with the same arithmetic exponent and polynomial bit overhead.

Now write an entry of the SAT coefficient ring in its polynomially many
integer module coordinates. Lift the coordinates from $\mathbb Z/2^p\mathbb Z$
to $[0,2^p)$. For each pair of module coordinates, compute the corresponding
integer matrix product using the rational implementation. The resulting
entries have bit length $O(p+\log d_2)$ for inner dimension $d_2$.
Combine them using the ring's integer multiplication table and reduce
modulo $2^p$. All ring dimensions, precisions, and logarithmic matrix
dimensions are polynomial in $s$. Coordinate pairs can be processed
sequentially. This proves the required bit-cost and storage bounds.

Alternatively, compute those integer products modulo enough good odd
primes and reconstruct them by CRT, specializing the fixed number field
through a fixed-degree quotient algebra. Only finitely many denominator
primes are excluded. No assertion about a dual exponent directly in
characteristic two is needed, and no nonunit is inverted in the SAT ring.

The fixed number-field schemes are existence ingredients of the
asymptotic algorithm. The delivered finite graph evaluator uses ordinary
exact multiplication; it does not implement these large fast schemes or
claim a practical crossover improvement.

\subsection{Operation supports}

For a merge of disjoint vertex sets $A,B$, let
$O=[r]\setminus(A\cup B)$. The three matrix axes are the independent
pair-index classes $A$--$B$, $A$--$O$, and $B$--$O$. The final triangle
is treated in the same way. Since the coefficients in either nontrivial
branch above are nonincreasing, the rearrangement inequality expresses
the sorted cost as the maximum over permutations of its coefficients.
Thus the classical integer supports are the permutations of
\[
 (56,56,0),\quad(56,47,30),
\]
with divisor $56$, and the additional backend uses
\[
 (44,44,0),\quad(44,35,20),
\]
with divisor $44$. There are nine distinct coefficient triples in each
case. An edge internal to one part gets coefficient zero; otherwise its
coefficient is that of its pair class. The support sets include every
binary merge and the final triangle of each plan.
